\section{Finite blocks determine weighted primitive Dirichlet factors}
\label{sec:prog}

In this section assume the critical-line statement for primitive Dirichlet
$L$-functions in the pinned source. For a primitive character $\chi$,
including the character of conductor one, let $m_\chi(\rho)$ be the multiplicity
of a nontrivial zero. Write
\[
   z_\rho=1-1/\rho,\qquad
   \lambda_\chi(n)=\lim_{R\to\infty}
      \sum_{|\rho|\le R}m_\chi(\rho)(1-z_\rho^n).
\]
Here and below sums using $m_\chi$ run over distinct zeros. These are the
starred Li sums, with the orientation of the Cayley map fixed as displayed.
The classical zero count gives $\sum m_\chi(\rho)|1-z_\rho|^2<\infty$,
and the starred first sum converges. See Lagarias \cite{LagariasLi}.
The only possible limit point of the $z_\rho$ on the circle is $1$.
The repository proves identification from all positive modes for primitive
products at one time, and identification in natural density for individual
characters. Classical comparison theorems can use the zero divisors
directly. For instance, Banks proves a short-height-window discrepancy for
a pair of distinct primitive characters of the same conductor
(\href{https://arxiv.org/abs/2302.07073}{arXiv:2302.07073}, Thm.~1.1).
Our data consist of exponentiated Li coefficients at finitely many indices
in a prescribed progression.

Fix a nonempty finite set $\mathcal C$ of distinct primitive characters,
a stride $h\ge1$, a residue $0\le r<h$, and a nonnegative integer $d$.
For $b\in\mathbb C^{\mathcal C}$ put
\[
          A_b(k)=\sum_{\chi\in\mathcal C} b_\chi\lambda_\chi(r+hk).
\]
When the weights are positive, $\exp(-A_b)$ is a product of
heat coefficients at the indicated times. Complex weights allow comparison
and subtraction without choosing a branch of a complex power. A calibration
polynomial $P$ multiplies the observed sequence by $\exp(P(k))$.
It includes an unknown normalization and rotation.

\begin{theorem}[Uniform separation on a finite block]\label{thm:calibrated}
For each $M>0$ there exist an integer $W\ge1$ and a constant $c>0$,
depending only on $\mathcal C,h,r,d,M$, such that every interval $I$
of $W$ consecutive positive integers and every $b$ with $\|b\|_2\le M$ satisfy
\begin{equation}
 \inf_{\substack{P\in\mathbb C[X]\\ \deg P\le d}}
     \max_{k\in I}\left|\exp\big(-A_b(k)+P(k)\big)-1\right|
      \ \ge\ c\,\|b\|_2 .
 \label{eq:calibration}
\end{equation}
The same constants apply when a different calibration polynomial is fitted
on each interval.
\end{theorem}

For a catalogue of $m$ characters, every admissible $W$ satisfies
$W\ge m+d+1$. Indeed, on $W$ points the values of the polynomials of
degree at most $d$ form a subspace of dimension $\min(d+1,W)$.
If its codimension is smaller than $m$, the linear map $b\mapsto A_b|_I$
followed by the quotient has a nonzero kernel element. Scale it to have
norm at most $M$ and interpolate its values to make the left side zero.

\begin{lemma}[Rigidity of a finite weighted zero discrepancy]\label{lem:logr}
If $\sum_{\chi\in\mathcal C} b_\chi m_\chi(\rho)=0$ at all but finitely
many nontrivial zeros in the union of the zero sets, then $b=0$.
This lemma holds without a critical-line hypothesis.
\end{lemma}

\begin{proof}
Let $\xi_\chi$ be the entire order-one completion of $L(s,\chi)$, including
the factor $s(s-1)$ for the conductor-one character. Its zeros are exactly
the nontrivial zeros. The genus-one Hadamard formula gives
\[
  \frac{\xi_\chi'}{\xi_\chi}(s)
   =B_\chi+\sum_\rho m_\chi(\rho)
     \left(\frac{1}{s-\rho}+\frac{1}{\rho}\right).
\]
The paired summands converge normally away from their poles. Thus the
$b$-weighted sum of these logarithmic derivatives is a rational function,
a constant plus finitely many simple fractions. Write
$V(s)=\sum_\chi b_\chi L'(s,\chi)/L(s,\chi)$. The logarithmic derivative
of a gamma factor is $\tfrac12\psi((s+a_\chi)/2)$ with
$a_\chi\in\{0,1\}$. Its increment under $s\mapsto s+2$ is
$1/(s+a_\chi)$, by the gamma recurrence. The conductor contributes a
constant and the elementary factors contribute rational functions.
Consequently $V(s+2)-V(s)$ is rational. On the positive real axis it is
$O(2^{-s})$ by absolute convergence of the logarithmic Euler series.
A rational function with this decay vanishes identically.
Iterating the resulting period to the right gives $V=0$ on $\Re s>1$.

Uniqueness of Dirichlet series now yields
$\sum_\chi b_\chi\chi(p)=0$ at every prime $p$ coprime to the least
common multiple $q$ of the conductors. Every reduced class modulo $q$
contains a prime. Distinct primitive characters induce distinct characters
of $(\mathbb Z/q\mathbb Z)^\times$, so orthogonality gives $b_\chi=0$
for all $\chi$.
\end{proof}

\begin{proof}[Proof of Theorem~\ref{thm:calibrated}]
Set $j=\max(2,d+1)$ and $\Delta f(k)=f(k+1)-f(k)$.
The sequence $u_b=\Delta^j A_b$ has an absolutely convergent expansion
\begin{equation}
 u_b(k)=-\sum_z B_z z^r (z^h-1)^j(z^h)^k,\qquad
 B_{z_\rho}=\sum_{\chi\in\mathcal C}b_\chi m_\chi(\rho).
 \label{eq:push-prog}
\end{equation}
Indeed $|z^h-1|\le h|1-z|$ and $j\ge2$.
It follows that $u_b$ is uniformly almost periodic on $\mathbb Z$ and
depends continuously, in the supremum norm, on $b$.

Suppose $\exp(-u_b(k))=1$ for all $k\in\mathbb Z$.
The values of $u_b$ then lie in $2\pi i\mathbb Z$.
An almost period with error less than $2\pi$ is an exact period. Choose
such a positive period $v$. By uniqueness of Fourier coefficients of
finite measures on the circle, the measure in \eqref{eq:push-prog} after
grouping at $w=z^h$ is supported on $w^v=1$.
Choose an open arc about $1$ on which $z\mapsto z^h$ is injective.
There are finitely many atoms outside the arc. Remove their images and
the $v$th roots of unity. Above every remaining $w$, any contributing
atom is the sole preimage in that arc. Its coefficient in
\eqref{eq:push-prog} must vanish, and $w\ne1$, so $B_z=0$.
The discarded finite set has finite fibers. Thus $B_z$ is zero except
at finitely many $z$. Lemma~\ref{lem:logr} implies $b=0$.
In particular $u_b$ itself is nonzero whenever $b\ne0$.

For $\|a\|_2=1$ and $0\le s\le M$ consider the bounded sequences
\[
 g_{a,s}=
 \begin{cases}
  (\exp(-s u_a)-1)/s,& s>0,\\
  -u_a,&s=0.
 \end{cases}
\]
They are uniformly almost periodic. The map $(a,s)\mapsto g_{a,s}$ is
continuous into the space of bounded sequences with its supremum norm,
including at $s=0$ by the exponential series. Each sequence is nonzero
by the preceding argument. For a nonzero uniformly almost-periodic
sequence choose an index where its modulus exceeds $2\eta>0$.
The translates of this index by its $\eta$-almost periods have bounded
gaps, so some excursion above $\eta$ occurs in every interval of a
fixed length. The same length works with threshold $\eta/2$ throughout a
supremum-norm neighborhood of the sequence. A finite subcover of the compact
parameter space gives $L$ and $\eta_0>0$ such that for every $b$ with
$0<\|b\|_2\le M$ there is a $k$ in each block of length $L$ with
\begin{equation}
          |\exp(-u_b(k))-1|\ge\eta_0\|b\|_2 .
 \label{eq:recurrentdiff}
\end{equation}

Write $R(k)=\exp(-A_b(k)+P(k))$. Since $\Delta^jP=0$,
\[
     \prod_{i=0}^{j} R(k+i)^{(-1)^{j-i}\binom{j}{i}}
       =\exp(-u_b(k)).
\]
This rational function of $j+1$ inputs equals $1$ at $(1,\ldots,1)$.
On the closed polydisc $|R_i-1|\le1/2$, its first derivatives are bounded,
so its difference from $1$ is at most $C_j\max_i|R_i-1|$.
Take $W=L+j$ and choose $k$ from its first $L$ positions using
\eqref{eq:recurrentdiff}. Either some $|R(k+i)-1|>1/2$, or the
last display gives $\max_i|R(k+i)-1|\ge (\eta_0/C_j)\|b\|_2$.
Taking $c=\min(1/(2M),\eta_0/C_j)$ proves \eqref{eq:calibration}.
\end{proof}

\paragraph{Reading the estimate.}
Suppose two weighted lists from the fixed catalogue give coefficient
sequences whose ratio, after calibration, differs from $1$ by at most
$\varepsilon$ on a block of length $W$. If their weight difference has
norm at most $M$, that difference has norm at most $\varepsilon/c$.
At zero error the weights agree exactly. At a common known time, this
recovers every factor multiplicity of a primitive product. If time is
unknown, the identifiable quantities are time times multiplicity:
$P$ at time $at$ and $P^a$ at time $t$ yield exactly the same data for every
positive integer $a$. A fixed roster and a weight bound allow a common
block length. The general proof gives existence of $W,c$. An explicit instance for two
real characters follows below.

For any two individual finite mixtures one can take the union of their
catalogues and choose $M$ after their weights are specified.
Consequently even a sequence of blocks whose union has natural density
zero determines a finite mixture, provided their lengths tend to infinity
and the ratio approaches $1$ in relative measure on them after
polynomial calibration of fixed degree. Indeed, for distinct mixtures
every full subblock of length $W$ contains an error of at least
$c\|b\|_2$, a positive lower fraction. For example, testing the modes
$r+hk$ with $2^{j^2}\le k<2^{j^2}+j$ suffices in this asymptotic sense.
This conclusion controls relative errors in the Fourier coefficients;
they can themselves be very small.

\subsection{An explicit 255-mode theorem for zeta and Dirichlet beta}
Let $\beta(s)=L(s,\chi_{-4})$ be the Dirichlet beta function. Its coefficients
and those of $\zeta$ are real, so their Li coefficients are real.
Write $A_{u,v}(k)=u\lambda_\zeta(r+20k)+v\lambda_\beta(r+20k)$ for
$0\le r<20$. The following constants apply to all starting indices.

\begin{theorem}\label{thm:explicit255}
For every interval $I$ of $255$ consecutive positive integers, all real
$u,v$, and every real affine function $P$, one has
\[
  \max_{k\in I}|A_{u,v}(k)-P(k)|
          \ \ge\ \tfrac16\max(|u|,|v|).
\]
Consequently the logarithms of the heat-coefficient magnitudes of
$\zeta^a\beta^b$ at time one, at these 255 indices in any residue class
modulo 20, determine the nonnegative integer multiplicities $a,b$ even
in the presence of unknown affine drift. Uniform errors strictly below
$1/12$ in the log magnitudes permit at most one such integer pair.
\end{theorem}

\begin{proof}
We give the numerical inputs and bounds explicitly. The first two positive
ordinates of $\zeta$ lie in $(14.13,14.14)$ and $(21,22)$.
The first four for $\beta$ lie in $(6.02,6.03)$, $(10.24,10.25)$,
$(12.98,13.00)$ and $(16.34,16.35)$. The certified data come from
\href{https://flintlib.org/doc/acb_dirichlet/zeta-zeros.html}{FLINT's
indexed zeta zero enclosures} and the
\href{https://www.lmfdb.org/L/1/2e2/4.3/r1/0/0}{LMFDB entry for $\chi_{-4}$}
with its linked zero list and reliability statement. Both spectra are
symmetric and have no central zero.

For $\gamma>0$ put $w(\gamma)=z_{1/2+i\gamma}^{20}
=\exp(40i\arctan(1/(2\gamma)))$.
Choose as targets $w_\zeta,w_\beta$ from the first positive zero of
each function. They satisfy, for either target $w_*$,
\begin{equation}
 |1-w_*|^2>5/3,\qquad |w-w_*|>3/25
       \quad\text{at every other folded zero in the two spectra}.
 \label{eq:boxes}
\end{equation}
Here is the full gap check. In increasing order of positive ordinate,
the real parts of the first four folded frequencies in the union lie in
\[
 (-.987,-.984),\quad (-.373,-.369),\quad
 (.030,.034),\quad (.155,.157).
\]
The fourth belongs to $\zeta$. All remaining positive ordinates give real
part greater than $.340$. Conjugate frequencies have the same real part.
The absolute imaginary parts at the two targets exceed $.16$ and $.98$,
respectively, separating each from its conjugate. For $\gamma\ge10$ the
angle in the exponential decreases in $(0,\pi)$, so its cosine increases.
On the first interval the angle lies between $\pi$ and $3\pi/2$.
Thus rational endpoint evaluations of
$((\gamma^2-1/4+i\gamma)/(\gamma^2+1/4))^{20}$ prove all these bounds.

We also need a bound on the mass of all the other frequencies. With
multiplicity in each sum, the two completed functional equations give
\begin{equation}
 \sum_{\rho\text{ of }\zeta}\frac1{|\rho|^2}
 +\sum_{\rho\text{ of }\beta}\frac1{|\rho|^2}
   =2-2\log\pi+2\frac{\beta'(1)}{\beta(1)}<\tfrac12 .
 \label{eq:bothmass}
\end{equation}
Indeed for either real completion $\xi$ the difference of its logarithmic
derivatives at $1$ and $0$ equals $2\xi'(1)/\xi(1)$ by symmetry.
The Hadamard expansion gives the same difference as
$\sum 1/(\rho(1-\rho))=\sum 1/|\rho|^2$ by the critical-line input.
For $\zeta$ the resulting expression is $2+\gamma_E-\log(4\pi)$, and for
$\beta$ it is $\log(4/\pi)-\gamma_E+2\beta'(1)/\beta(1)$, where
$\gamma_E$ is Euler's constant. To check the strict bound, use
$\beta(1)=\pi/4>0.78$, $\log\pi>1.14$, and $\beta'(1)<0.30$.
For the last inequality the alternating series for $\beta'(1)$ is bounded
above by the five terms from $\log(3)/3$ through $\log(11)/11$.
Termwise differentiation at $1$ follows from uniform convergence of the
alternating derivative series on a real neighborhood of $1$.
The script \texttt{check\_window.py} verifies these elementary
inequalities and the frequency enclosures using rational arithmetic.

Write $D(k)=A_{u,v}(k+1)-2A_{u,v}(k)+A_{u,v}(k-1)$ and
$B=\max(|u|,|v|)$. Its Fourier series has an atom for each zero at
$w=z^{20}$ with coefficient $b_\rho z^r|1-w|^2$, where $b_\rho$ is the
corresponding scalar $u$ or $v$, counted with the multiplicity of the zero.
By $|1-z^{20}|\le20|1-z|$ and \eqref{eq:bothmass}, the sum of the absolute
coefficients is less than $200B$ when $B>0$. Choose the target belonging to
a scalar of modulus $B$. It has coefficient of absolute value greater
than $5B/3$ by \eqref{eq:boxes}.

Let $K(X)=(64^{-1}\sum_{j=0}^{63}X^j)^4=\sum_{j=0}^{252}a_jX^j$.
Its coefficients are nonnegative and sum to one. In the average
$\sum_{j=0}^{252}a_j\overline w_*^{\,j}D(k_0+j)$ the chosen atom contributes
its full coefficient times a unit factor. Every other frequency is
multiplied in modulus by at most $(25/96)^4$, by the geometric sum and
\eqref{eq:boxes}. Since $200(25/96)^4<1$, the average has absolute
value greater than $2B/3$. Some $|D(k_0+j)|$ therefore exceeds $2B/3$.
An affine function has zero second difference, and the absolute coefficient
sum of a second difference is $4$. At least one of the $255$ indices
from $k_0-1$ to $k_0+253$ has $|A_{u,v}-P|>B/6$.
The case $B=0$ is immediate. For two integer pairs, their nonzero difference
has $B\ge1$. Two fits with errors strictly below $1/12$ would violate the
inequality by the triangle inequality.
\end{proof}
